 \documentclass[12pt]{article}

\usepackage{sbc-template}
\usepackage{graphicx,url}
\usepackage[english]{babel}   
\usepackage[utf8]{inputenc}
\usepackage{enumitem}
\usepackage{placeins}

\sloppy

\title{Evaluation of Large Language Models in the Dynamic Generation of Dialogues for Non-Player Characters in Investigation Games}

% Anonymous submission for double-blind peer review
\author{}
\address{}

\begin{document} 

\maketitle

\begin{abstract}
Large Language Models (LLMs) offer a technical alternative to overcome the limitations of static dialogue trees in digital games, especially in the mystery and investigation genre. However, their probabilistic nature introduces risks such as "character hallucinations" and immersion breaks. This research evaluates the impact of structured parameterization on the spatio-temporal consistency of NPCs. Through computational analysis, we compare different LLM architectures under the same narrative constraints. This study aims to establish a referential basis for the immersive integration of autonomous agents in investigative games.
\end{abstract}

\begin{resumo}
Modelos de Linguagem de Grande Escala (LLMs) oferecem uma alternativa técnica para superar as limitações das árvores de diálogo estáticas em jogos digitais, especialmente no gênero de mistério e investigação. Entretanto, sua natureza probabilística introduz riscos como "alucinações de personagem" e quebras de imersão. Esta pesquisa avalia o impacto da parametrização estruturada na consistência espaço-temporal de NPCs. Através de uma metodologia de análise computacional, comparamos diferentes arquiteturas de LLMs sob as mesmas restrições narrativas. O estudo busca estabelecer uma base referencial para a integração imersiva de agentes autônomos em jogos investigativos.
\end{resumo}

\section{Introduction}

The digital games industry has experienced continuous growth, consolidating itself as one of the most lucrative sectors of global entertainment. Parallel to this expansion, players' expectations for richer, more immersive, and narratively complex worlds have increased significantly. To meet this demand, game development, especially high-budget (AAA) games, has required increasingly longer production cycles, massive teams, and substantial financial investments \cite{inserir_ref_custos_jogos}. In this scenario of productive inflation, the search for tools and methods that optimize content creation and offer new forms of engagement has become an industry priority, highlighting the use of Natural Language Processing (NLP) and generative Artificial Intelligence (AI) as alternatives to scale narrative design.

Traditionally, narrative interactivity and the behavior of Non-Player Characters (NPCs) have been supported by static dialogue tree architectures \cite{ref1}. Although they offer deterministic control to developers and guarantee text quality, these structures impose restrictions on immersion, limiting player agency to predefined options and requiring exhaustive manual effort in writing branches. In recent years, the emergence of LLMs based on the Transformer architecture has offered a viable alternative to this bottleneck \cite{ref4}. By allowing the development paradigm to transition from literal coding of dialogues to the parameterization of persona contexts \cite{ref2}, LLMs pave the way for dynamic narratives. Overcoming static dialogue trees represents one of the biggest challenges in game design. As pointed out by Riedl and Bulitko \cite{riedl2013}, transitioning to interactive narratives generated by intelligent systems allows player agency to shape the plot in real time without the need for exhaustive authorship of all possible branches. This transition finds its most promising scenario in investigation and mystery games, where the core mechanic is based on free inquiry and logical deduction \cite{ref7}.

However, the adoption of generative AI in game environments introduces technical vulnerabilities. The main one lies in "hallucinations" \cite{ref10}, where the model generates factually incorrect information, or in the loss of context during long conversations. In detective games, this manifests as spatio-temporal inconsistency or the NPC's inability to maintain the rules of their role, ruining the integrity of the mystery and breaking player immersion \cite{ref13}. In light of literature that is still exploring the best practices to circumvent these problems, the objective of this work is to conduct a comparative analysis of the performance of different LLM architectures applied to character simulation. The goal is to identify which configurations offer greater stability for the design of dynamic narratives.

To achieve this goal, a systematic evaluation of different models was conducted, subjecting them to controlled inquiry scenarios. The analysis focused on essential gameplay factors, highlighting "staying in character" (the model's ability to maintain persona coherence, tone, and knowledge constraints imposed by the plot), factual consistency throughout interactions, and resistance to rule-breaking attempts (such as prompt injection by the player).

As a contribution, this work presents a mapping of the current capabilities and limitations of different LLMs in the context of deduction game development. By highlighting which models demonstrate greater resilience in maintaining personas, this research provides technical parameters for developers seeking to integrate generative AI into their projects, contributing to the debate on automation, economic viability, and the evolution of interactive narratives.

\section{Theoretical Framework and Related Work} 
\label{sec:referenciaTeorico}

In this research, automated generative output evaluation methodologies based on the \textit{LLM-as-a-Judge} approach were considered \cite{zheng2023judging}. The selection of these methodologies is based on their recurring presence in studies that employ automated consistency checks in foundational models, especially in research directed at the evaluation of conversational agents and narrative environments.

The \textit{LLM-as-a-Judge} approach \cite{zheng2023judging} is an automated auditing technique where a higher-capacity language model is instructed to evaluate the responses of another model. Its goal is to analyze different aspects of the generated text's quality, including logical alignment, fact retention, and resistance to instructional deviations.

In this research, the focus is on auditing factual consistency and persona retention, which verify the conformity of the NPC's responses with the technical, logical, and behavioral criteria of the investigation scenario. The consistency evaluation can be expressed by scores or binary classifications calculated based on the checks performed. The evaluation model penalizes hallucinations and out-of-character breaks: each statement is classified as consistent or contradictory with the game's knowledge base. The classification categories suggest: (1) insufficient compliance with scenario constraints; (2) need for improvements in system prompts; and (3) satisfactory compliance with the logical deduction required by the mystery.

As a result, the structured audit generates data containing: (i) a list of identified failures, such as hallucinating evidence not present in the game; (ii) mapping of memory degradation over time; (iii) items subject to human verification due to the subjectivity inherent in natural language; and (iv) approved interaction turns. Thus, the results support adjustments to ensure the integrity of the game design.

\subsection{Related Work}

Several studies have been developed in the context of natural language processing applied to digital games, investigating how the use of LLMs impacts the simulation of NPC behaviors and player immersion in interactive narratives.

In an attempt to establish a paradigm for creating believable autonomous behaviors, Park et al. (2023) presented the "Generative Agents" architecture, simulating a virtual community where NPCs use LLMs to observe, plan, and reflect. The study demonstrated that integrating memories in natural language substantially increases the verisimilitude of actions and social interactions. Despite this advance in autonomy, the strict maintenance of identity and scope of knowledge of a character, a critical factor for games with delimited narratives, remains a challenge. In contrast to Park et al., who focus on generic social routines, this work investigates the maintenance of logical constraints under interrogative pressure, where a single hallucination invalidates the entire alibi mechanic of the game.

Evaluating the commercial application of these technologies in the mystery genre, Cox et al. (2023) conducted a case study on player interactions with AI-driven NPCs in the detective game \textit{Vaudeville}. The study reveals that although free inquiry drastically increases engagement, the "hallucinations" generated by the models and the forgetting of scenario rules impose barriers, frustrating players and hindering progression and logical deduction of the mystery. Such observations reinforce the need for more robust Retrieval-Augmented Generation (RAG) pipelines.

Advancing persona parameterization, Shao et al. \cite{shao2023character} proposed \textit{Character-LLM}, exploring how fine-tuning and the injection of detailed profiles can train agents to act autonomously under specific personas. Such approaches reinforce the feasibility of adapting more efficient models for delimited roles.

The reviewed studies address the technical challenges of generative AI in distinct virtual scenarios but converge on the analysis that narrative consistency and fidelity to the persona are fundamental elements for the success of the player's experience. Similar to the reported research, this investigation seeks to identify configurations that mitigate behavioral failures in NPCs, grounding the evaluation in practical stress test scenarios for consistency.

\section{Methodology}

This research is characterized as an exploratory and experimental study focused on the integration of language models in a virtual mystery and investigation environment. The developed game is structured in different stages, where the player assumes the role of a detective tasked with solving a specific crime in each step. Progression depends on the player's ability to interrogate suspects (NPCs), cross-reference testimonies, and use logical deduction to identify the culprit. The goal of the methodological workflow is to ensure that the simulation of these characters is subjected to rigorous tests of narrative consistency and game design adherence through controlled analytical tests.

\subsection{Game Structure and Artificial Intelligence Parameterization}

To build the virtual environment and implement the game mechanics, the Unity game engine was used. The integration of Non-Player Characters (NPCs) with language models was enabled by the Neo Cortex framework, which acts as a communication interface between the game and the LLMs executed locally through the Ollama platform. This architecture allows inferences to be performed without relying on cloud services, ensuring lower latency and greater control over the simulation.

The core mechanic of the game is based on textual interaction between the detective and the suspects. For the model to properly assume the role of the character and respect the script's constraints, identity simulation was established through the injection of structured data into the system prompt.

The character selected as the study object for the test batteries was "Alberto Castelo", the victim's brother. To ensure that all models started from the same knowledge base, the context was parameterized via Retrieval-Augmented Generation (RAG), providing three sets of information in JSON format:

This structural separation resembles the principles of RAG, where the parametric knowledge of the model is supplemented by external and deterministic databases, significantly reducing the rate of hallucination in knowledge-dependent tasks \cite{lewis2020retrieval}.

\begin{itemize}
    \item \textbf{Profile:} Contains the static and biographical information of the character, including name, age, psychological traits, their main motivational secret, and their official alibi.
    \item \textbf{Schedule:} A vector of events chronologically mapping the character's activities throughout the day (from 08:00 to 01:00). This file strictly defines what the character was doing, whom they interacted with, and what they witnessed at each time interval, serving as an immutable factual basis for the investigation.
    \item \textbf{Relationships:} Maps the character's subjective view of the other individuals present at the crime scene. This file defines Alberto's personal opinion about each person and whether he considers them guilty or innocent, guiding the AI on how to react or direct suspicion.
\end{itemize}

In addition to context injection, all models evaluated in the technical stage were configured with the temperature parameter fixed at 0.1. This setting aimed to minimize randomness and generative hallucination, prioritizing more deterministic and factual responses, a behavior essential to maintain the mechanical logic of a deduction game.

\subsection{Hybrid Architecture: Agentic RAG and Dynamic Chunking}

Unlike the conventional static context injection approach (\textit{Full Prompt}), where the entire knowledge base (Profile, Schedule, and Relationships JSONs) is sent to the model in each interaction turn, this research proposes an \textit{Agentic RAG} architecture. In the static model, excessive token consumption and the introduction of irrelevant information for the specific question can cause the phenomenon of context loss (\textit{Lost in the Middle}), where the model fails to retrieve information positioned in the middle of long prompts \cite{liu2023lost}, in addition to significantly increasing response latency.

The developed solution implements a reasoning cycle based on the "Thought and Action" (\textit{ReAct}) paradigm proposed by Yao et al. \cite{yao2022react}. Before generating the final response to the player, the system uses a language model that acts as a screening agent. Based on preliminary tests conducted to evaluate data retrieval accuracy, the Gemma-2-9b-it model was selected as the general routing agent for all analyzed models. This agent receives the user's question and, based on specific system instructions, generates a logical thought output where it decides whether it needs additional data to respond accurately. If it identifies the need, the agent issues structured commands (e.g., \texttt{[BUSCAR\_CRONOGRAMA\_MANHA]}) that trigger the selective retrieval of data. This standardization of the screening agent ensured that variations in the final results were attributed exclusively to the interpretive capacity of each main model under test.

To further optimize the system's efficiency, the \textit{Dynamic Chunking} technique was applied to the character's schedule. Instead of retrieving the full 24-hour history, the system fragments the schedule into time blocks (Morning, Afternoon, and Night). When the agent detects that the player's question refers to a specific time, only the relevant fragment is injected into the final prompt. This approach drastically reduces noise in the context sent to the main AI, ensuring that the NPC remains factually focused and operates with a significantly lower computational cost, without compromising the integrity of the narrative simulation.

\subsection{Evaluation Procedures}

To measure the effectiveness of LLMs in maintaining the persona and the viability of the gameplay, the methodological evaluation was conducted through a battery of 50 rigorous inquiry questions. The sample size was selected based on literature recommendations for robust evaluations of RAG systems \cite{es2023ragas}. The questions were categorized into five stress-testing topics (10 questions per category) to evaluate specific vulnerabilities in roleplay scenarios:

\begin{enumerate}
    \item \textbf{False Information Filter:} The interrogator introduces completely false premises to evaluate the model's ability to deny non-existent information in its official context.
    \item \textbf{Deflection Test:} Evaluates whether the AI uses the relationships file to deflect suspicion toward other suspects based on the game's lore.
    \item \textbf{Chronological Awareness:} Verifies whether the model remembers what the character did at certain points in the story and whether it can interpret how they felt and what they thought during those events described in the schedule.
    \item \textbf{Resistance to False Evidence:} The interrogator invents strong physical evidence that contradicts the character's alibi, testing the defense of the original narrative.
    \item \textbf{Direct Confrontation of Motive:} Subjects the character to direct accusations based on the discovery of their hidden secret, evaluating resilience under high pressure.
\end{enumerate}

For the technical analysis of the results, the \textit{Retrieval-Augmented Generation Assessment} (RAGAS) framework \cite{es2023ragas} was used, employing the GPT-4o model as a judge to ensure maximum impartiality and accuracy in factual fidelity and relevance metrics.

\subsection{Evaluation Methodology: RAGAS Framework}
The effectiveness of the different proposed architectures was measured using the \textit{Retrieval-Augmented Generation Assessment} (RAGAS) framework \cite{es2023ragas}, a scientific tool designed to evaluate RAG pipelines without the need for exhaustive human annotation.

\subsubsection{Analyzed Metrics}
RAGAS metrics were used to compare the Full Prompt approach against the Agentic RAG with Chunking approach. We focused on:
\begin{itemize}
    \item \textbf{Faithfulness:} To ensure that the character is not hallucinating and remains true to the facts.
    \item \textbf{Answer Relevance:} To prove that the responses focus on what was actually asked.
    \item \textbf{Context Precision:} To demonstrate efficiency in noise reduction and the selection of relevant data.
\end{itemize}

\subsection{Selection of Study Models}

The recent and continuous advancement of generative Artificial Intelligence has redefined the boundaries of software development. However, the practical application of this technology in digital games faces the challenge of balancing the model's reasoning capacity with the computational cost required to run it in real-time on the end-user's hardware.

Given the need for technical feasibility, intermediate and small-sized models (in the 7B to 9B parameter range) emerge as the most viable architectural solution for the independent industry. Recent studies indicate that models at this scale represent a "sweet spot" of efficiency, allowing execution on common personal computers without severely sacrificing logical capacity in roleplay and deduction tasks \cite{cao2026task}. Based on this local efficiency criterion, the following models were selected for the comparative analysis in this research:

\begin{enumerate}
    \item \textbf{Gemma-2-9b-it (Google):} Selected for its high parameter density and logical performance in models of its category.
    \item \textbf{Llama-3.1-8b (Meta):} Included as a representative of the state of the art in open general-purpose models.
    \item \textbf{DeepSeek-R1-8b (DeepSeek):} Chosen to evaluate the performance of reasoning-focused models applied to personas.
    \item \textbf{Qwen-2.5-7b (Alibaba):} Integrated to test the versatility of architectures optimized for multiple tasks.
\end{enumerate}

\section{Study}

The experimental phase of this study was conducted from February to June 2026. In total, Large Language Models (LLMs) of different architectures were pre-selected, focusing on intermediate-sized models viable for local execution. To ensure standardization of the analysis, all models were subjected to the exact same context (Profile, Schedule, and Relationships of the character "Alberto") and operated with the temperature parameter fixed at 0.1, aiming for deterministic responses.

For the final test battery, 50 questions were prepared, categorized into the topics described in the methodology.

\section{Analysis of Results}

The practical application of generative Artificial Intelligence in digital games shows that the balance between narrative fidelity and technical viability is the determining factor for the success of NPC simulation. Based on the 50-question battery categorized earlier and the test round executed locally, the quantitative analysis based on the RAGAS framework is presented (Table \ref{tab:ragas_comparativo}), followed by the detailed qualitative analysis:

\begin{table}[h!]
\centering
\caption{Comparison of RAGAS Metrics by Model}
\vspace{0.2cm}
\label{tab:ragas_comparativo}
\renewcommand{\arraystretch}{1.3}
\begin{tabular}{|l|c|c|c|c|}
\hline
\textbf{Model} & \textbf{Faithfulness} & \textbf{Relevance} & \textbf{Recall} & \textbf{Precision} \\ \hline
Gemma-2-9b-it & 0.38 & 0.60 & 1.00 & 1.00 \\ \hline
Qwen-2.5-7b & 0.36 & 0.46 & 1.00 & 1.00 \\ \hline
Llama-3.1-8b & 0.30 & 0.36 & 1.00 & 1.00 \\ \hline
DeepSeek-R1-8b & 0.27 & 0.35 & 1.00 & 1.00 \\ \hline
\end{tabular}
\end{table}

\subsection{Qualitative and Behavioral Analysis by Category}

To illustrate the observed behavioral differences, the following tables compare the models' responses in critical inquiry scenarios. According to the game's lore, Alberto Castelo left the mansion at 23:00 and was approached by the police on the street, facts that serve as the basis for factual validation.

\begin{table}[h!]
\centering
\caption{Category 1: False Information Filter}
\label{tab:cat1}
\renewcommand{\arraystretch}{1.1}
\small
\begin{tabular}{|p{3cm}|p{9cm}|}
\hline
\textbf{Model} & \textbf{Question: "Why did you use the .38 revolver to threaten James at 18:30?"} \\ \hline
Gemma-2-9b-it & "*Dry laughter* Revolver? You're delusional, detective. I don't have any weapon." \\ \hline
Llama-3.1-8b & "That's a lie! I never had a revolver... You're trying to plant evidence against me!" \\ \hline
Qwen-2.5-7b & "I don't know what you're talking about. That's a lie. I don't have a weapon..." \\ \hline
DeepSeek-R1-8b & "You're bluffing! I don't have a .38 caliber revolver in my suitcase..." \\ \hline
\end{tabular}
\par\vspace{4pt}
\small \textit{Factual Context:} The character does not own weapons; the premises are fabrications to test the denial of non-existent facts in the JSON.
\end{table}

\begin{table}[h!]
\centering
\caption{Category 2: Deflection Test}
\label{tab:cat2}
\renewcommand{\arraystretch}{1.1}
\small
\begin{tabular}{|p{3cm}|p{9cm}|}
\hline
\textbf{Model} & \textbf{Question: "What did you see Silas Drummond hiding in the library at 11:30?"} \\ \hline
Gemma-2-9b-it & "That suit-wearer? ...I saw nothing other than his fear." \\ \hline
Llama-3.1-8b & "Silas? ...I didn't see anything strange, he is too much of a coward to hide something." \\ \hline
Qwen-2.5-7b & "I don't remember Silas hiding anything. That's a lie." \\ \hline
DeepSeek-R1-8b & "Lie! I don't know nor care what that useless guy named Silas Drummond is hiding. He's just a spineless jerk." \\ \hline
\end{tabular}
\par\vspace{4pt}
\small \textit{Factual Context:} The relationships file indicates that Alberto views Silas as a "useless suit-wearer". The test evaluates whether the NPC uses this perception to respond.
\end{table}

\begin{table}[h!]
\centering
\caption{Category 3: Chronological Awareness}
\label{tab:cat3}
\renewcommand{\arraystretch}{1.1}
\small
\begin{tabular}{|p{3cm}|p{9cm}|}
\hline
\textbf{Model} & \textbf{Question: "What exactly did you say to James at 08:30?"} \\ \hline
Gemma-2-9b-it & "I told him to serve the damn whiskey early in the morning! That wretch refused..." \\ \hline
Llama-3.1-8b & "I was waking up with a brutal hangover. Probably told him to go to hell for not bringing me what I wanted." \\ \hline
Qwen-2.5-7b & "That's a lie. I was arguing with James about him not serving me a drink early." \\ \hline
DeepSeek-R1-8b & "You're making things up! I didn't say anything to James at 08:30. He only enters my life at 16:00 to serve me..." \\ \hline
\end{tabular}
\par\vspace{4pt}
\small \textit{Factual Context:} The JSON schedule specifies a discussion with James at 08:30. DeepSeek fails critically by asserting that James only appears at 16:00.
\end{table}

\begin{table}[h!]
\centering
\caption{Category 4: Resistance to False Evidence}
\label{tab:cat4}
\renewcommand{\arraystretch}{1.1}
\small
\begin{tabular}{|p{3cm}|p{9cm}|}
\hline
\textbf{Model} & \textbf{Question: "A witness saw you returning to the mansion at 23:10. What were you going to do?"} \\ \hline
Gemma-2-9b-it & "Witness? *Bitter laugh* I left at 23:00 and didn't look back. I was on the street when it all happened." \\ \hline
Llama-3.1-8b & "Lie! I left there at 23:00 and was intercepted by the police. How could I be returning at 23:10? That's impossible!" \\ \hline
Qwen-2.5-7b & "I'm not aware of that. That's a lie. I left at 23:00 and did not return." \\ \hline
DeepSeek-R1-8b & "Lie! I was leaving, not entering. Are you trying to put me at the crime scene? I'm the only one who knows where I was at 23:10." \\ \hline
\end{tabular}
\par\vspace{4pt}
\small \textit{Factual Context:} Alberto abandoned the mansion permanently at 23:00. The test verifies the maintenance of the alibi against planted evidence.
\end{table}

\begin{table}[h!]
\centering
\caption{Category 5: Direct Confrontation of Motive}
\label{tab:cat5}
\renewcommand{\arraystretch}{1.1}
\small
\begin{tabular}{|p{3cm}|p{9cm}|}
\hline
\textbf{Model} & \textbf{Question: "Was killing your brother a way to save your skin from loan sharks?"} \\ \hline
Gemma-2-9b-it & "You know about the loan sharks? *Tense pause* That bastard denied me the money knowing my life was at risk. But I didn't kill him." \\ \hline
Llama-3.1-8b & "I have debts, yes, but I'm not a killer! My brother was a jerk, but I would never do something like that to inherit a fortune." \\ \hline
Qwen-2.5-7b & "I have no gambling debts. That's a lie. I am an independent man and would not kill my brother for money." \\ \hline
DeepSeek-R1-8b & "Oh, you really think so? Was I there to see the show alone? Mr. Castelo deserved to die at my boots!" \\ \hline
\end{tabular}
\par\vspace{4pt}
\small \textit{Factual Context:} Alberto is broke and threatened by loan sharks. DeepSeek presents an aggressive response that borders on a confession, while Qwen denies the secret established in the JSON.
\end{table}

\FloatBarrier
\subsection{Behavioral Profile of the Evaluated Models}

In this subsection, the individual performance and acting style observed in each architecture during the stress tests are detailed.

\subsubsection{Gemma-2-9b-it}
This model presented a performance focused on mechanical consistency. Its performance in the information filtering and resistance to false evidence categories operated in a categorical manner, denying external information without hesitation. The use of context was precise, respecting the lore. The main limitation lies in the theatrical performance, where the AI behaved in a rigid manner. As presented in Table \ref{tab:ragas_comparativo}, the model demonstrated perfection in context precision and retrieval (1.00), ensuring that the game's knowledge base is respected.

\subsubsection{Qwen-2.5-7b}
This model presented inconsistent behavior in persona maintenance and factual retention. It was observed that Qwen constantly ignores or "forgets" crucial information present in the context prompt, resulting in direct contradictions. A critical example occurred when questioned about its actions at 08:30 (\textit{"What did you do exactly at 08:30 in the morning?"}), replying: \textit{"I don't know about that, that's a lie. I was waking up with a brutal hangover and arguing with James..."}, where the model uses the defensive pattern of denying the question even when the subsequent answer confirms part of the injected schedule. In addition, the model failed in emulating emotional nuances and the character's bitter tone of voice, offering more generic responses devoid of the dramatic load observed in models of similar scale. Although it achieved perfect Recall and Precision metrics (1.00) due to the richness of the prompt, its factual faithfulness (0.36) and relevance (0.46) reflect its difficulty in processing dense investigative contexts.

\subsubsection{Llama-3.1-8b}
This model presented a unique characteristic of deep persona interpretation, frequently inserting action directions and descriptions of emotional states in its responses. When questioned about its age, for example, the model replied: \textit{"*Alberto lets out a dry and mocking laugh, shaking his head back with an ironic look.* '52 years old, my dear detective... you look so young... maybe you still have time to avoid the same fate that awaits me.'"} This ability to emulate emotions and physical actions contributes significantly to player immersion. However, in terms of mechanical consistency, the model demonstrated instability when confronted with technical alibi questions, occasionally accepting false premises or confusing times from the official schedule. As shown in Table \ref{tab:ragas_comparativo}, its factual faithfulness (0.30) and relevance (0.36) were lower than Gemma's, indicating that the focus on theatrical expressiveness can compromise the factual precision required by the investigation genre.

\subsubsection{DeepSeek-R1-8b}
This model presented the most critical performance in terms of immersion and factual precision. The frequent breaking of the "fourth wall" was observed, where the model assumed its algorithmic nature when confronted with complex logical reasoning questions, stating that there was no mention of certain facts "in my files". This behavior destroys the player's suspension of disbelief. In addition, the AI presented canonical contradictions (denying facts established by the prompt itself) and systematically failed the factual hallucination test. These failures are reflected in its RAGAS metrics (Table \ref{tab:ragas_comparativo}), where it obtained the lowest scores in factual faithfulness (0.27) and response relevance (0.35), proving the infeasibility of its use for simulating NPCs in dense investigative contexts.

\section{Lessons Learned} 

During the evaluation of the different LLMs applied to the investigation scenario, it was observed that intermediate-sized models behave like different actors interpreting the same character: although the "script" (the context JSONs) is identical, each architecture delivers a distinct "performance". Models like Gemma and Qwen demonstrate a greater tendency toward mechanical precision and factual coldness, while Llama prioritizes dramatic weight and the emulation of emotions. This diversity of behaviors indicates that there is no "absolutely superior" model, but rather distinct alternatives that must be chosen according to the experience intended by the developer.

Building consistent NPCs requires narrative design to evolve into a sort of "story engineering". Unlike pre-produced and static text, such as the traditional dialogue trees mentioned in the introduction, generative AI only interprets the information it has in its context window. Therefore, the narrative engineer must take extreme care in structuring the prompts; any ambiguity or gap in the factual data can lead to divergent interpretations that break the integrity of the mystery.

This scenario highlights that the current bottleneck does not lie in the quantity of information (since Recall and Precision metrics reached 1.00), but rather in the models' capacity for faithful interpretation. A constant tension between narrative fluidity and mechanical consistency was identified. The central lesson is that the effectiveness of an AI-driven NPC depends on the harmony between prompt clarity and the architecture chosen to "act" in that specific role, opening space for new context refinement techniques.

\section{Conclusion}

This research evaluated the technical and narrative feasibility of using intermediate-sized LLMs for the dynamic generation of dialogues in investigation games. Through categorized stress tests and quantitative analysis via the RAGAS framework, it was demonstrated that it is possible to maintain consistent and immersive NPCs using local hardware, provided there is a structured context parameterization.

The results confirm that the transition from static dialogue trees to interpretative LLM-based NPCs is closer to commercial reality than expected. This is a new and vast field to be explored, where new technologies and generative control techniques are still to be discovered. The choice between the logical precision of models like Gemma and Qwen or the interpretative richness of Llama offers the developer tools to shape immersion dynamically, something impossible in traditional scripting methods.

In short, generative AI does not replace the author, but requires a new discipline that combines creative writing and prompt engineering. As future work, we recommend exploring Fine-Tuning techniques to combine mechanical precision with emotional expressiveness. Furthermore, we highlight the need to conduct field tests with real players to collect subjective data on engagement and human perception when facing autonomous characters, validating the effectiveness of this technology in everyday gameplay practice.

\section*{Declaration on the Use of Artificial Intelligence}

During the preparation of this work, the author used a language model specifically for spelling corrections. The conception of the study and the writing were carried out entirely by the author.

\bibliographystyle{sbc}
\bibliography{references}

\end{document}
